% ---------------------------------------------------------------------------
%  02-theory.tex -- the theorem
% ---------------------------------------------------------------------------
\section{The decomposition theorem}
\label{sec:theory}

Write $\atoms{\varphi}$ for the atomic propositions occurring in $\varphi$ and
$\clocks{\varphi}$ for the formula-clocks of the timed modalities occurring in
it. A node's clocks are allocated per modality, so
$\atoms{\varphi}\cap\atoms{\psi}=\emptyset$ for every pair drawn from different
components already gives $\clocks{\cdot}$ disjoint as well; the two are not
independent assumptions.

\begin{definition}[Split]\label{def:split}
A node $N=\tup{\Gam}{\Ctx}{\zone}{\obligations}{\Deadline}$ with pursuit set
$P$ \emph{splits} into $N_1,\dots,N_r$ ($r\ge 2$) when
\begin{enumerate}[label=\textup{(S\arabic*)},ref=\textup{(S\arabic*)}]
\item\label{S1} $\Gam$, $\Ctx$ and $P$ are partitioned among the $N_i$, and
      $\atoms{N_i}\cap\atoms{N_j}=\emptyset$ for $i\neq j$;
\item\label{S2} every clock constrained in $\zone$ belongs to exactly one
      $\clocks{N_i}$;
\end{enumerate}
and $N_i$ carries $\proj{\zone}{\clocks{N_i}}$, the deadlines on its own clocks,
and the phase of $N$.
\end{definition}

There is deliberately no condition on the zone. An earlier version required
$\zone$ to be the product of its projections,
$\zone=\bigcap_i\proj{\zone}{\clocks{N_i}}$, on the grounds that a reset leaves
finite difference constraints in the DBM and can relate $x$ and $y$ even when
the modalities owning them share no atom. That condition is sufficient but
\emph{not} necessary, and the proof below never uses it: see
\cref{rem:noproduct}. \ref{S2} is a real requirement, and it is about a clock
owned by no group at all: a modality still sitting in $\Gam$, not yet launched, already
owns its clock, and that clock may already carry a value.

\begin{theorem}[Decomposition]\label{thm:decomp}
If $N$ splits into $N_1,\dots,N_r$, then $N$ is realisable if and only if every
$N_i$ is.
\end{theorem}

\begin{proof}
($\Rightarrow$) Let $\mathcal T$ realise $N$, with location $s$ and the
interpretation of the formula-clocks fixed by $\obligations$. Fix $i$ and let
$\mathcal T_i$ be $\mathcal T$ with its labelling restricted to $\atoms{N_i}$
and its clocks restricted to $\clocks{N_i}$. Every formula of $N_i$ mentions
only atoms in $\atoms{N_i}$ and is evaluated against clocks in
$\clocks{N_i}$, and the satisfaction relation of TCTL inspects a configuration
only through those two, so truth of each such formula is unchanged by the
restriction; the same runs remain, with the same delays, so
$\Runsinf$ is unchanged and the quantifiers range over the same set. Every
$\nu\in\proj{\zone}{\clocks{N_i}}$ extends, by definition of the projection, to
some $\nu'\in\zone$, and $\mathcal T,(s,\nu')$ satisfies the formulas of $N_i$, hence so does
$\mathcal T_i,(s,\nu)$. Thus $\mathcal T_i$ realises $N_i$.

($\Leftarrow$) Let $\mathcal T_i$ realise $N_i$ at $(s_i,\nu_i)$, induced by a
TA $\mathcal A_i$ over clock set $\clocks{N_i}$. The clock sets are pairwise
disjoint, so form the product $\mathcal A=\prod_i\mathcal A_i$ in the standard
way: locations are tuples, the invariant of a tuple is the conjunction of the
components' invariants, the labelling is the union of the components'
labellings --- consistent, because the atom sets are disjoint --- and the edges
interleave, each edge of one $\mathcal A_i$ lifting to an edge of the product
that leaves the other components' locations and clocks untouched. Delays
synchronise: a delay of $t$ in the product is a delay of $t$ in every component,
admissible because each component's invariant permits it.

Three things now hold of $\mathcal A$.

\emph{Projection of runs.} Every run $\rho$ of $\mathcal A$ projects to a run
$\proj{\rho}{i}$ of $\mathcal A_i$ with the same delays, and $\rho$ is
non-Zeno iff its projections are, the delays being shared. Conversely, given a
run $\rho_i$ of $\mathcal A_i$ and any runs of the others, the delays can be
merged into a single run of the product projecting to each: refine the
components' delay sequences to their common refinement and interleave the
discrete steps in any order consistent with time. Non-blockingness of each
component gives non-blockingness of the product, so $\Runsinf$ is non-empty
everywhere.

\emph{Truth is componentwise.} By induction on $\varphi$ with
$\atoms{\varphi}\subseteq\atoms{N_i}$: for a literal this is the definition of
the product labelling; the boolean cases are immediate; for
$\mathrm{E}(\cdot)$ take the witnessing run of $\mathcal A_i$ and lift it
against arbitrary runs of the others, and for $\mathrm{A}(\cdot)$ take any run
of the product and project it to $\mathcal A_i$, where the hypothesis applies.
The timed operators add only a constraint on $\mathit{time}(\pi)$, which the
projection preserves because delays are shared.

\emph{The zone.} Let $\nu\in\zone$. By \ref{S2} each of its coordinates belongs
to exactly one component, so $\nu$ decomposes as $(\nu_1,\dots,\nu_r)$ with
$\nu_i\in\proj{\zone}{\clocks{N_i}}$, and $(s_i,\nu_i)$ realises $N_i$'s
formulas. By the previous paragraph the product configuration
$((s_1,\dots,s_r),\nu)$ satisfies every formula of every $N_i$, that is, all of
$\Gam\cup\Ctx$. As $\nu\in\zone$ was arbitrary, $\mathcal A$ realises $N$.
\end{proof}

\begin{remark}[No zone-product condition is needed]\label{rem:noproduct}
Take $\zone$ constrained by $x-y<1$ with $x,y$ in different components: both
projections are universal and $\zone$ is not, so $\zone$ is a proper subset of
the product of its projections. Splitting is still sound, and each direction
says why. Forwards, a valuation of a projection extends to one of $\zone$, and a
component's formulas read only its own clocks. Backwards, each component is
asked to be realisable over its \emph{projection}, which is the larger demand:
$\zone\subseteq\bigcap_i\proj{\zone}{\clocks{N_i}}$, so a product model meeting
the projections meets $\zone$ a fortiori.

Dropping the condition is a simplification of the proof rather than a
measured gain: on every instance in the suite it never once blocked a split
that \ref{S1} and \ref{S2} allowed, and node counts are identical with it and
without it.
\end{remark}

\begin{corollary}[What the search may do]\label{cor:search}
If $N$ splits, its verdict in the tableau is the conjunction of the verdicts of
its components.
\end{corollary}

\begin{proof}
The calculus is sound and complete, so a node lies in the greatest successful
marking exactly when it is realisable; apply \cref{thm:decomp}.
\end{proof}

\begin{intuition}
Two halves of a node that never mention the same symbol and never share a clock
are two separate questions, even if the zone happens to relate their clocks ---
neither half can read the other's clock, so the relation is invisible to both.
Answer the two questions separately and combine the answers by running the two
models side by side: they agree about the passage of time and about nothing
else, which is exactly what a product automaton is.
\end{intuition}
